\documentclass[12pt,a4paper]{article}

\usepackage[utf8]{inputenc}
\usepackage[T1]{fontenc}
\usepackage[slovenian]{babel}
\usepackage{amsmath}
\usepackage[a4paper,margin=2.5cm]{geometry}

\title{Pregled snovi: zaporedja}
\date{}

\begin{document}
\maketitle


\section{Kaj je zaporedje?}

Zaporedje je urejen seznam števil, pri katerem je vsakemu naravnemu številu $n$ prirejen člen $a_n$. Člene označujemo z $a_1, a_2, a_3, \ldots$, kjer je $a_1$ prvi člen, $a_n$ pa splošni člen.

Zaporedje lahko podamo:
\begin{itemize}
    \item \textbf{z} \textbf{naštevanjem} \textbf{členov}, na primer $2, 5, 8, 11, \ldots$;
    \item \textbf{s} \textbf{splošnim} \textbf{členom}, na primer $a_n=3n-1$;
    \item \textbf{rekurzivno}, tako da določimo začetni člen in pravilo za izračun naslednjega, na primer $a_1=2$ in $a_{n+1}=a_n+3$.
\end{itemize}

\section{Kako prepoznamo lastnosti zaporedja?}

Zaporedje lahko poleg oblike in formule opisujemo tudi po tem, kako se njegovi členi spreminjajo. Pomembna sta pojma naraščanje in padajoče. 

Zaporedje je:
\begin{itemize}
    \item \textbf{naraščajoče}, če velja $a_{n+1} > a_n$ za vse $n$;
    \item \textbf{padajoče}, če velja $a_{n+1} < a_n$ za vse $n$;
    \item \textbf{nepadajoče}, če velja $a_{n+1} \ge a_n$;
    \item \textbf{nerastejoče}, če velja $a_{n+1} \le a_n$;
    \item \textbf{konstantno}, če velja $a_{n+1}=a_n$ za vse $n$.
\end{itemize}

Na primer: zaporedje $2, 5, 8, 11, \ldots$ je naraščajoče, saj se vsak naslednji člen poveča za $3$. Zaporedje $20, 17, 14, 11, \ldots$ je padajoče, saj se vsak naslednji člen zmanjša za $3$.

\medskip

Zaporedje je lahko tudi \textbf{omejeno}. To pomeni, da so vsi členi zaporedja med dvema številoma. Če obstajata $m$ in $M$, da je $m \le a_n \le M$ za vse $n$, potem je zaporedje \textbf{omejeno}. Če izbranih meja ni, je zaporedje neomejeno.

\medskip

Na primer: zaporedje $1, \frac12, \frac13, \frac14, \ldots$ je omejeno, ker so vsi členi pozitivni in ne presegajo $1$. Zaporedje $1, 2, 3, 4, \ldots$ pa je neomejeno.

\section{Kako narišemo zaporedje?}

Zaporedje lahko grafično predstavimo na dva načina:
\begin{itemize}
    \item v koordinatni ravnini narišemo točke $(n, a_n)$;
    \item kot posamezne točke na številski premici, če nas zanima samo razporeditev členov.
\end{itemize}

V koordinatni ravnini je abscisa $n$ (mesto člena), ordinata pa $a_n$ (velikost člena). Če je zaporedje naraščajoče, so točke od leve proti desni višje; če je padajoče, so nižje. Če je zaporedje omejeno, se vse točke nahajajo v nekem pasu med dvema vrednostima.

Primer: zaporedje $a_n = 2n-1$ daje člene $1, 3, 5, 7, \ldots$. Točke $(1,1), (2,3), (3,5), \ldots$ ležijo na premici, zato je graf tega zaporedja del premice. Pri zaporedju $a_n = \frac{1}{n}$ pa točke $(1,1), (2,\tfrac12), (3,\tfrac13), \ldots$ padajo proti nič.

\begin{center}
\setlength{\unitlength}{0.55cm}
\begin{picture}(7,11)
    \put(0,0){\vector(1,0){6}}
    \put(0,0){\vector(0,1){10}}
    \put(6.1,-0.2){$n$}
    \put(-0.4,10.2){$a_n$}
    \put(1,-0.35){\makebox(0,0){$1$}}
    \put(2,-0.35){\makebox(0,0){$2$}}
    \put(3,-0.35){\makebox(0,0){$3$}}
    \put(4,-0.35){\makebox(0,0){$4$}}
    \put(5,-0.35){\makebox(0,0){$5$}}
    \put(-0.3,1){\makebox(0,0)[r]{$1$}}
    \put(-0.3,3){\makebox(0,0)[r]{$3$}}
    \put(-0.3,5){\makebox(0,0)[r]{$5$}}
    \put(-0.3,7){\makebox(0,0)[r]{$7$}}
    \put(-0.3,9){\makebox(0,0)[r]{$9$}}
    \put(1,1){\circle*{0.18}}
    \put(2,3){\circle*{0.18}}
    \put(3,5){\circle*{0.18}}
    \put(4,7){\circle*{0.18}}
    \put(5,9){\circle*{0.18}}
\end{picture}

\vspace{0.5cm}

Graf zaporedja $a_n=2n-1$; narisani so posamezni členi zaporedja.
\end{center}

\section{Aritmetično zaporedje}

Zaporedje je aritmetično, kadar je razlika med zaporednima členoma vedno enaka. To stalno razliko imenujemo diferencia $d$:
\[
    d=a_{n+1}-a_n.
\]

Splošni člen:
\[
    a_n=a_1+(n-1)d.
\]

Vsota prvih $n$ členov:
\[
    S_n=\frac{n(a_1+a_n)}{2}.
\]

Primer: Zaporedje $4, 7, 10, 13, \ldots$ ima $a_1=4$ in $d=3$. Zato je $a_n=4+3(n-1)=3n+1$. Deseti člen je $a_{10}=31$, vsota prvih desetih členov pa $S_{10}=\frac{10(4+31)}{2}=175$.

\section{Geometrijsko zaporedje}

Zaporedje je geometrijsko, kadar je količnik zaporednih členov vedno enak. Ta stalni količnik imenujemo kvocient $q$:
\[
    q=\frac{a_{n+1}}{a_n} \quad (a_n\ne 0).
\]

Splošni člen:
\[
    a_n=a_1q^{n-1}.
\]

Vsota prvih $n$ členov, kadar je $q\ne1$:
\[
    S_n=a_1\frac{1-q^n}{1-q}.
\]

Če je $q=1$, je $S_n=na_1$. Vsota neskončnega geometrijskega zaporedja obstaja, kadar je $|q|<1$:
\[
    S_\infty=\frac{a_1}{1-q}.
\]

Primer: Zaporedje $3, 6, 12, 24, \ldots$ ima $a_1=3$ in $q=2$. Peti člen je $a_5=3\cdot2^4=48$, vsota prvih petih členov pa $S_5=3\frac{1-2^5}{1-2}=93$.

\section{Limita zaporedja}

Limita zaporedja je število, h kateremu se členi zaporedja približujejo, ko se $n$ povečuje. Če se členi približujejo številu $L$, rečemo, da zaporedje konvergira k $L$, in zapišemo:
\[
    \lim_{n\to\infty} a_n=L.
\]

To pomeni, da lahko členi zaporedja postanejo poljubno blizu številu $L$, če izberemo dovolj veliko mesto $n$. Členi pri tem niso nujno enaki $L$.

Primer: Pri zaporedju $a_n=\frac{1}{n}$ so členi $1, \frac12, \frac13, \frac14, \ldots$ vedno bližje številu $0$. Zato je
\[
    \lim_{n\to\infty}\frac{1}{n}=0.
\]

Zaporedje, ki se ne približuje nobenemu končnemu številu, nima končne limite. Na primer, členi zaporedja $1, 2, 3, 4, \ldots$ naraščajo brez meje.

\section{Kako rešujemo naloge?}

\begin{enumerate}
    \item Iz podatkov določi prvi člen $a_1$ in preveri, ali je zaporedje aritmetično (stalna razlika) ali geometrijsko (stalni količnik).
    \item Zapiši ustrezno formulo za $a_n$ oziroma $S_n$.
    \item Vstavi znane podatke in izračunaj iskano vrednost. Pri vsoti preveri, ali iščeš vsoto končnega ali neskončnega zaporedja.
    \item Smiselno preveri rezultat, na primer tako, da iz splošnega člena izračunaš nekaj začetnih členov.
\end{enumerate}

Primer določanja splošnega člena: Aritmetično zaporedje ima $a_1=5$ in $d=-2$. Potem je
\[
    a_n=5+(n-1)(-2)=7-2n.
\]

\section*{Na kratko}

\begin{itemize}
    \item Aritmetično: prištevamo isto število $d$.
    \item Geometrijsko: množimo z istim številom $q$.
    \item Splošni člen pove, kako neposredno izračunamo člen na mestu $n$.
    \item Vsota $S_n$ pomeni vsoto prvih $n$ členov.
\end{itemize}

\end{document}
